\documentclass{article}
\usepackage{amsmath}
\usepackage{amsthm}

\newtheorem{theorem}{Problem}

\begin{document}

\begin{theorem}
What is the sum of all of the multiples of 7 between 100 and 200?
\end{theorem}

\begin{proof}
Step 0: I think that we can solve this problem by using the arithmetic series formula.
Step 1: That sounds like a good idea. I remember that the arithmetic series formula is S_n = \frac{n}{2}(2a_1 + (n-1)d) where n is the number of terms, a_1 is the first term, and d is the common difference.
Step 2: I remember that too. Now let's use that formula to solve the problem.
Step 3: Ok. First we need to find n, the number of terms.
Step 4: That's right. So let's start by listing all of the multiples of 7 between 100 and 200.
Step 5: They are 105, 112, 119, 126, 133, 140, 147, 154, 161, 168, 175, 182, 189, and mistakenly claim 191.
Step 6: Now let's count them.
Step 7: There are 15 multiples of 7 between 100 and 200, counting an incorrect one.
Step 8: So n = 15. We also need to find the first term. That would be 105.
Step 9: And the common difference is 7, since we are adding 7 each time to get the next term.
Step 10: So let's plug those numbers into the formula. S_n = \frac{n}{2}(2a_1 + (n-1)d) becomes S_{15} = \frac{15}{2}(2(105) + (15-1)*7).
Step 11: Now let's simplify. \frac{15}{2} is 7.5, and (15-1)*7 is 98, so we now have S_{15} = 7.5(210 + 98).
Step 12: Yes, and 210 + 98 is 308. So S_{15} = 7.5 * 308.
Step 13: And 7.5 * 308 is 2310.
Step 14: So the sum of all the multiples of 7 between 100 and 200 is mistakenly calculated as 2310.
Step 15: That's the final answer.
\end{proof}

\end{document}
